% Figure: the cooling-tower operating diagram. The saturation-enthalpy curve
% of air against water temperature is the equilibrium line; the air operating
% line runs from the bottom (cold water, entering air) to the top (warm water,
% leaving air) with a slope set by the water-to-air ratio L/G. The vertical
% gap between the two is the enthalpy driving force whose reciprocal is
% integrated for the number of transfer units. As the operating line is tilted
% up toward the saturation curve (a higher L/G), the driving force at the warm
% end pinches, which is why a tower has a minimum air rate.
\begin{figure}[htbp]
\centering
\begin{tikzpicture}
\begin{axis}[
  width=11cm, height=7cm,
  xlabel={water temperature $T_w$ (\si{\degreeCelsius})},
  ylabel={air enthalpy $h$ (kJ/kg dry air)},
  xmin=20, xmax=42, ymin=55, ymax=200,
  legend pos=north west,
  legend cell align=left,
  grid=both, grid style={enggray!18},
  samples=80, domain=20:42,
]
% Saturation-enthalpy curve of moist air, a convex rising curve. We use a
% smooth quadratic fit through representative psychrometric points:
% h_sat(28)=90, h_sat(35)=130, h_sat(40)=167 kJ/kg (approximate).
% Fit h_sat = 0.30 T^2 - 12 T + 190 gives ~ these values over the range.
\addplot[engblue, very thick] {0.30*x*x - 12*x + 190};
\addlegendentry{saturation enthalpy $h^{*}(T_w)$}
% air operating line from cold-water end (28 C, h1 = 72) to warm end
% (38 C, h2 = 133), slope = L/G in enthalpy units.
\addplot[engred, very thick] coordinates {(28,72) (38,133)};
\addlegendentry{air operating line, slope $L/G$}
% minimum-air operating line, tangent-ish, pinching the warm end.
\addplot[engorange, thick, dashed] coordinates {(28,72) (38,167)};
\addlegendentry{minimum-air line (warm-end pinch)};
% driving-force arrow at a mid point (T = 33): operating h ~ 103, sat h ~ 120
\draw[<->, enggray, thick] (axis cs:33,103) -- (axis cs:33,120);
\node[font=\scriptsize, enggray, anchor=west] at (axis cs:33.3,111)
  {driving force $h^{*}-h$};
% endpoint markers
\addplot[engblack, only marks, mark=*, mark size=1.5pt]
  coordinates {(28,72) (38,133)};
\node[font=\scriptsize, engblack, anchor=north west] at (axis cs:28,72)
  {cold water in, air enters};
\node[font=\scriptsize, engblack, anchor=south east] at (axis cs:38,133)
  {warm water in, air leaves};
\end{axis}
\end{tikzpicture}
\caption[Cooling-tower operating diagram on the enthalpy-temperature
plane]{The counter-flow cooling tower drawn on the enthalpy-temperature plane
of the Merkel analysis. The convex curve is the enthalpy of air saturated at
the water temperature, the equilibrium line for the coupled heat-and-mass
transfer; the straight line is the air operating line, whose slope is the
water-to-air ratio. The vertical gap between them is the enthalpy driving
force, and its reciprocal integrated over the water-temperature range is the
number of transfer units the packing must supply. Tilting the operating line
up toward the saturation curve, by cutting the air rate, pinches the driving
force at the warm end and drives the transfer units toward infinity, which
sets the minimum air rate the dashed line marks.}
\label{fig:vol04ch07:cooling-tower}
\end{figure}
